\section*{Chapter \ref{ch:g1:what-is-a-computer} --- What Is a Computer?}

\begin{solution}{exo:g1:what-is-a-computer:1}
The phone, the oven that stops by itself and the car contain a \omterm{def:g1:what-is-a-computer:computer}{computer}. The
lamp with a plain switch and the wooden chair follow no instructions: no
\omterm{def:g1:what-is-a-computer:computer}{computer}.
\end{solution}

\begin{solution}{exo:g1:what-is-a-computer:2}
\omterm{def:g1:what-is-a-computer:device}{Input devices}: the keyboard, the mouse, the microphone, the camera. \omterm{def:g1:what-is-a-computer:device}{Output
devices}: the screen, the printer, the speaker, the headphones.
\end{solution}

\begin{solution}{exo:g1:what-is-a-computer:3}
\omterm{def:g1:what-is-a-computer:input}{Input}: the picture that the camera takes (and Lina's touch on the button).
\omterm{def:g1:what-is-a-computer:input}{Output}: the photo shown on the screen.
\end{solution}

\begin{solution}{exo:g1:what-is-a-computer:4}
It shows pictures to us (\omterm{def:g1:what-is-a-computer:input}{output}), and it tells the \omterm{def:g1:what-is-a-computer:computer}{computer} where our finger
touches it (\omterm{def:g1:what-is-a-computer:input}{input}).
\end{solution}

\begin{solution}{exo:g1:what-is-a-computer:5}
Desktop \omterm{def:g1:what-is-a-computer:computer}{computer}: the screen, the keyboard, the mouse. Laptop: two of the
screen, the camera, the touchpad (and its keyboard).
\end{solution}

\begin{solution}{exo:g1:what-is-a-computer:6}
For example: the \omterm{def:g1:what-is-a-computer:computer}{computers} of the \omterm{def:g1:what-is-a-computer:computer}{computer} room, the teacher's laptop, the
tablets of the class, the photocopier, the \omterm{def:g1:what-is-a-computer:computer}{computer} inside the school bell or
the heating.
\end{solution}

\begin{solution}{exo:g1:what-is-a-computer:7}
No. Many \omterm{def:g1:what-is-a-computer:computer}{computers} have no screen: the \omterm{def:g1:what-is-a-computer:computer}{computer} inside a washing machine, a
car key or a toy.
\end{solution}

\begin{solution}{exo:g1:what-is-a-computer:8}
\omterm{def:g1:what-is-a-computer:input}{Input}: Sami's spoken words, through the microphone (an \omterm{def:g1:what-is-a-computer:device}{input device}). \omterm{def:g1:what-is-a-computer:input}{Output}:
the song, through the speaker (an \omterm{def:g1:what-is-a-computer:device}{output device}).
\end{solution}

\begin{solution}{exo:g1:what-is-a-computer:9}
Yes. Its \omterm{def:g1:what-is-a-computer:input}{input} is the keys you press; it follows instructions to work out the
answer; its \omterm{def:g1:what-is-a-computer:input}{output} is the answer on its screen.
\end{solution}

\begin{solution}{exo:g1:what-is-a-computer:10}
\omterm{def:g1:what-is-a-computer:input}{Inputs}: the turn of the dial, the level of the water in the drum. \omterm{def:g1:what-is-a-computer:input}{Outputs}: the
motor that turns the drum, the tap that lets water in, the beep at the end.
\end{solution}

\begin{solution}{exo:g1:what-is-a-computer:11}
In: ``the robot touches a wall'' (and the button that starts it). Out: the
motors that turn its wheels, which stop when it touches the wall.
\end{solution}
